\section{Manga109Script}
Manga109Script is a script dataset corresponding to 19,555 manga pages from the Manga109 dataset~\cite{matsui2017sketch,aizawa_building_2020}. Manga109 consists of 109 commercial Japanese manga works published between the 1970s and 2010s, ensuring substantial variety in style and content. Manga109Script reformats the original two-page spreads into single pages and links each page to its corresponding script data.

The script format follows a style commonly used in Japanese manga production. Each script consists of three main components:
(1) Scene Headings, marked with a leading “◯” to denote location,
(2) Action Descriptions, written with a three-character indent to describe characters’ actions and environmental details, and
(3) Dialogue, where the speaker’s name is followed by their speech enclosed in quotation marks.
These components are illustrated in the center of Figure~\ref{fig:eyecatch}.

We adopted this script format because manga has primarily developed within Japan, and all manga titles in this study were originally created through Japanese production processes. This ensures the dataset aligns with actual manga creation practices.

\subsection{Data Construction Process}
To construct Manga109Script, we employed GPT-4o via API to generate page-level scripts from manga images in the Manga109 dataset. To support the model's understanding of character identity and dialogue context, we introduced two key inputs: (1) visual prompting, where the face of each character in the image was enclosed with a uniquely colored circle~\cite{shtedritski_what_2023}, and (2) pre-extracted dialogue annotations, including the textual content of each speech balloon and its associated speaker's name~\cite{li_manga109dialog_2024}. Dialogue order was pre-determined using the Magi algorithm~\cite{sachdeva_manga_2024}, which provides a reliable rule-based ordering of panels and speech balloons. Without this step, GPT-4.1 alone achieved exact ordering on only 26 out of 109 sampled pages, whereas the Magi-based ordering reached 96 out of 109 pages, demonstrating the necessity of algorithmic ordering prior to script generation.

To further address the challenge of recognizing continuous narratives across multiple pages, we structured a unique structured interaction process with GPT-4o consisting of the following three components for each target page~$i$:

\begin{itemize}
    \item \textbf{$\alpha(i)$: Initial comprehension prompting} \\
    We provided the full dialogue from the episode that includes page~$i$ and simulated an affirmative response from the model to promote narrative understanding.

    \vspace{0.5em}
    \begin{adjustwidth}{1em}{0pt}
      \setlength{\tabcolsep}{2pt}
      {\ttfamily\spaceskip=0.4em plus 0.1em minus 0.1em
      \begin{tabular}{@{}r l@{}}
      \textbf{Input:} & entire dialogue of the episode\\
      \textbf{Task:} & comprehend the episode \\
      \textbf{Agent:} & comprehension complete \\
      \end{tabular}
      }
  \end{adjustwidth}
  \vspace{0.5em}

    \item \textbf{$\beta$($i$): Script generation for the $i$-th page} \\
    We provided GPT-4o with the processed image of page~$i$ (with visual prompting), the dialogue on page~$i$, and the summary generated in step $\gamma(i{-}1)$, covering the story up to page~$i{-}1$. The model was then instructed to generate the script for page~$i$ in the Japanese screenwriting format.

    \vspace{0.5em}
    \begin{adjustwidth}{0.3em}{0pt}
      \setlength{\tabcolsep}{2pt}
      {\ttfamily\spaceskip=0.5em plus 0.1em minus 0.1em
      \begin{tabular}{@{}r l@{}}
      \textbf{Input:} & page 
  ~$i$ image and dialogue,\\
                      & and prior summary from $\gamma(i{-}1)$\\
      \textbf{Task:} & generate script for page~$i$  \\
      \textbf{Agent:} & script generation complete \\
      \end{tabular}
      }
  \end{adjustwidth}
  \vspace{0.5em} 

    \item \textbf{$\gamma$($i$): Summarization up to the $i$-th page} \\
    We then prompted the model to summarize the story up to page~$i$ to generate the summary input for $\beta(i+1)$.

    \vspace{0.5em}
    \begin{adjustwidth}{0.3em}{0pt}
    \setlength{\tabcolsep}{2pt}
    {\ttfamily\spaceskip=0.4em plus 0.1em minus 0.1em
    \begin{tabular}{@{}r l@{}}
    \textbf{Task:} & based on interaction history, \\
                   & summarize the story so far  \\
    \textbf{Agent:} & summarization complete \\
    \end{tabular}
    }
    \end{adjustwidth}
\end{itemize}
\vspace{0.5em} 

To generate the script for page~$i$, we executed the sequence:
\[
\alpha(i) \rightarrow \beta(i{-}1) \rightarrow \beta(i).
\]
Here, the arrow ($\rightarrow$) indicates the chronological order in which each prompt was issued to GPT-4o.

To generate the summary of the story up to page~$i$, we extended the sequence as follows:
\[
\alpha(i) \rightarrow \beta(i{-}1) \rightarrow \beta(i) \rightarrow \gamma(i).
\]

Note that for $i = 1$, these sequences become $\alpha(1) \rightarrow \beta(1)$ and $\alpha(1) \rightarrow \beta(1)\rightarrow \gamma(1)$, where $\beta(1)$ does not include a summary input.

By iteratively applying this series of interactions from $i = 1$ to the final page of each manga title, we constructed the full Manga109Script dataset. The dataset statistics are summarized in Table~\ref{tab:statistics}.

\begin{table}[t]
\setlength{\abovecaptionskip}{2pt}
\caption{
Statistics of Manga109Script. Distinct character types are counted by character identity per page, and distinct speaker types are counted by speaker identity per page.
}
{\setlength{\tabcolsep}{5pt}
\centering
{
\begin{tabular}{l r r}
\toprule
\textbf{Element} & \textbf{Count} & \textbf{Avg./page} \\
\midrule
\multicolumn{3}{l}{\textit{Layout Elements}} \\
\hspace{0.5em}Pages & 19{,}555 & \textemdash\\
\hspace{0.5em}Panels & 103{,}850 & 5.31 \\
\hspace{0.5em}Character instances & 157{,}865 & 8.70 \\
\hspace{0.5em}Distinct character types & 60{,}812 & 3.11 \\
\hspace{0.5em}Dialogue balloons & 146{,}549 & 7.49 \\
\hspace{0.5em}Distinct speaker types & 55{,}778 & 2.85 \\
\midrule
\multicolumn{3}{l}{\textit{Script Elements}} \\
\hspace{0.5em}Scene headings & 26{,}931 & 1.38 \\
\hspace{0.5em}Action descriptions & 138{,}890 & 7.04 \\
\hspace{0.5em}Dialogues & 134{,}901 & 7.10 \\
\hspace{0.5em}Japanese characters (text) & 6{,}703{,}063 & 342.78 \\
\bottomrule
\end{tabular}
}}
\label{tab:statistics}
\vspace{-1.2em}
\end{table}

\subsection{Quality Verification}
%The generated scripts have a specialized format unlike short captions. Hence, standard image caption evaluation metrics were difficult to apply. 
Unlike general image captioning datasets, our scripts lack ground truth.
Moreover, the script format provides contextual coherence rather than word overlap.
Thus, we used human evaluation with specific metrics to verify script quality.

Evaluation items included: (1) {\it Scene Heading Recall}, determining whether locations depicted in the manga were correctly identified in script scene headings; (2) {\it Scene Heading Precision}, determining whether locations mentioned in script scene headings actually appeared in the manga; (3) {\it Action Description Precision}, determining whether action descriptions generated for a target page actually match that page when mixed with randomly selected action descriptions from other pages; (4) {\it Out-of-Context Action Description Detection Accuracy}, determining whether action descriptions from other pages could be correctly identified as not belonging to the target page; and (5) {\it Dialogue Recall}, determining whether dialogue appearing in the manga was accurately transcribed in the script.

We randomly selected approximately 10 consecutive pages from each work to maintain story continuity, totaling 1,000 pages. The 1,000 pages were divided into two sets of 500 pages each, with six evaluators each reviewing one set, resulting in three evaluations per page. These evaluators received fair compensation.

We calculated each evaluation metric by averaging correct judgments across all elements: Scene Heading Recall (0.793), Scene Heading Precision (0.838), Action Description Precision (0.664), Out-of-Context Action Description Detection Accuracy (0.910), and Dialogue Recall (0.988).